\documentclass[10pt,a4paper]{amsart}
\usepackage[margin=19mm]{geometry}
\usepackage{amsmath,amssymb,mathtools}
\usepackage[scr=rsfs,cal=euler]{mathalfa}
\usepackage{xfrac}
\usepackage[unicode,hidelinks,bookmarks=false]{hyperref}
\newcommand{\ie}{i.e.\ }
\newcommand{\cf}{cf.\ }
\newcommand{\resp}{respectively\ }
\newcommand{\Subsection}{\subsection}
\providecommand{\nobreakdash}{\nobreak\hbox{-}}
\setlength{\emergencystretch}{2em}
\multlinegap0pt
\usepackage{xcolor}
\usepackage{amssymb}
\usepackage{mathtools}
\usepackage{tcolorbox}
\usepackage{cancel}
\usepackage[shortlabels]{enumitem}
\usepackage{bm}
\newtheorem{lemma}{Lemma}
\newcommand{\dif}{{\mathrm{d}}}
\newcommand{\mathbbm}[1]{\boldsymbol{#1}}
\title{Simple Approximation and Derivative Free Inference-Time Scaling for Diffusion Models via Sequential Monte Carlo on Path Measures}
\author{Chenyang Wang, Weizhong Wang, Yinuo Ren, Jose Blanchet, Yiping Lu}
\date{Source: arXiv:2605.17850v1 (CC BY 4.0)}
\begin{document}
\maketitle

\noindent\emph{Source and licence.} Simple Approximation and Derivative Free Inference-Time Scaling for Diffusion Models via Sequential Monte Carlo on Path Measures. Version arXiv:2605.17850v1 (CC BY 4.0), distributed by arXiv under the Creative Commons Attribution 4.0 International licence, http://creativecommons.org/licenses/by/4.0/, as recorded in the arXiv licence metadata for this exact version and shown on the abstract page https://arxiv.org/abs/2605.17850v1. Full licence notice: \texttt{licenses/arxiv-2605.17850-CC-BY-4.0.txt}.\par\smallskip

\noindent\textit{Editorial scope.} Counted unit: the article's lemma lem:convergence\_of\_discrete\_time\_generator (source0.tex lines 934--940). The article's complete original proof of it is delivered with it (source0.tex lines 943--1032, one \texttt{proof} environment of 90 lines). No mathematics is written, repaired or re-derived: the statement and the proof are the article's own lines, in the article's own notation, and no retained line is substituted or re-flowed.  No further source-local context is needed: every cross-reference of every retained line resolves inside the two delivered spans.  Nothing was omitted from any retained span, so the package carries no omission record.  No retained line contains a citation, so the wrapper carries no bibliography and no bibliography entry.  No retained line contains a figure environment, an image-inclusion command or drawing code, so no image is delivered and the package declares no asset dependency.  Editorial formatting only, in addition: an AMS \texttt{amsart} preamble that restates the article's own declarations and macros used by the retained lines, namely the environments \texttt{lemma} on separate counters, together with the standard packages those lines need.  The retained environments print in this document as Lemma 1 in order of appearance, whereas the article prints the same items with the numbering of its own full document.  The complete native source of the article as distributed in the arXiv e-print for this exact version is bundled unchanged as \texttt{sources/200820/source0.tex}.
\begin{lemma}
The discrete-time generator $L^{\Delta t}_t$ converges uniformly to the generator $\mathcal{L}_t$ of the limiting continuous-time jump-diffusion on the core $C_c^2(E)$, defined as 
\begin{equation*}
    \mathcal{L}_t f(\mathbbm{z}) = \mathcal{L}^{\text{diff}}_t f(\mathbbm{z}) + \Lambda(\mu^N_{t-}) \int_{[0,1]^N} \left( f(\Psi(\mathbbm{z}, \mathbbm{v})) - f(\mathbbm{z}) \right)  \mathrm{d}\mathbbm{v}.
\end{equation*}
\label{lem:convergence_of_discrete_time_generator}
\end{lemma}

\begin{proof}[Proof of Lemma~\ref{lem:convergence_of_discrete_time_generator}]
    We prove the lemma in two steps.
    
    \textbf{Step 1: Generator of the Continuous Process.}
    By It\^o's formula,
    \begin{align*}
    \dif r(X_t^G,t)
    &= \partial_t r(X_t^G,t)\dif t +\nabla_x r(X_t^G,t)^{\top}\big(v(X_t^G,t)+V^2(t)\nabla_x G(X_t^G,t)\big)\dif t\\
    &\quad + \frac12 V^2(t)\Delta_x r(X_t^G,t)\dif t
    + V(t)\nabla_x r(X_t^G,t)^{\top} \dif W_t.
    \end{align*}
    We use a generalized version of dynamics of the continuous process $\mathbbm{z}_t$:
    \[
    d\mathbbm{z}_t = v_{\text{diff}}(\mathbbm{z}_t) \dif t + \sigma_{\text{diff}}(\mathbbm{z}_t) \dif W_t + \int_{0}^{\infty} \int_{[0,1]^N} (\Psi(\mathbbm{z}_{t-}, \mathbbm{v}) - \mathbbm{z}_{t-}) \mathbbm{1}_{\{u \le \Lambda(\mu^N_{t-})\}} \mathcal{N}(\dif t, \dif u, \dif \mathbbm{v}),
    \]
    where 
    \[
        \begin{aligned}
            v_{\text{diff}}(\mathbbm{z}_t,t) =& \bigg(v(X^{(i),G}_t,t)+V^2(t)\nabla_x G(X^{(i),G}_t,t), \partial_t r(X_t^G,t)+\nabla_x r(X_t^{(i),G},t)^{\top}\big(v(X_t^{(i),G},t)+V^2(t)\nabla_x G(X_t^{(i),G},t)\big)\\
            &+\frac12 V^2(t)\Delta_x r(X_t^{(i),G},t)-\frac{1}{2}\|\theta(X^{(i),G}_t,t)\|^2\bigg)_{i=1}^N,
        \end{aligned}
    \]
    and
    \[
        \sigma_{\text{diff}}(\mathbbm{z}_t,t) = \left(V(t), V(t)\nabla_x r(X_t^{(i),G},t)^{\top}-V(t) \nabla_x G(X^{(i),G}_t,t)^\top\right)_{i=1}^N.
    \]
    By applying It\^o's formula for semimartingales with jumps, the time evolution of $f(\mathbbm{z}_t)$ is given by:
    \begin{align*}
    f(\mathbbm{z}_t) - f(\mathbbm{z}_0) &= \int_0^t \mathcal{L}^{\text{diff}}_t f(\mathbbm{z}_s) ds + \int_0^t \nabla f^\top \sigma_{\text{diff}} \dif W_s \\
    &\quad + \int_0^t \int_0^\infty \int_{[0,1]^N} \left( f(\Psi(\mathbbm{z}_{s-}, \mathbbm{v})) - f(\mathbbm{z}_{s-}) \right) \mathbbm{1}_{\{u \le \Lambda(\mu^N_{s-})\}} \mathcal{N}(\dif s, \dif u, \dif \mathbbm{v}),
    \end{align*}
    where the generator for the diffusion part is defined as
    \begin{equation*}
        \mathcal L_t^{\text{diff}} f(\mathbbm{z}) = v_{\text{diff}}(\mathbbm{z}, t)^\top \nabla f(\mathbbm{z}) + \frac{1}{2} \mathrm{tr}(\sigma^2_{\text{diff}}(\mathbbm{z}, t)\nabla^2 f(\mathbbm{z})).
    \end{equation*}
    Applying It\^o's formula for semimartingales to a test function $f \in C_c^2(E)$, and introducing the compensated Poisson measure 
    $$
        \widetilde{\mathcal{N}}(\dif s,\dif u, \dif \mathbbm{v}) = \mathcal{N}(\dif s,\dif u, \dif \mathbbm{v}) - \dif s \dif u \dif \mathbbm{v},
    $$ 
    we decompose the evolution of $f(\mathbbm{z}_t)$ into drift and martingale components. Since integrals with respect to $W_t$ and $\widetilde{\mathcal{N}}$ are local martingales with vanishing expectations, the generator $\mathcal{L}$ is defined by the predictable drift terms:
    \[
    \mathcal{L}_tf(\mathbbm{z}) = \mathcal{L}^{\text{diff}}_tf(\mathbbm{z}) + \int_0^\infty \int_{[0,1]^N} \left( f(\Psi(\mathbbm{z}, \mathbbm{v})) - f(\mathbbm{z}) \right) \mathbbm{1}_{\{u \le \Lambda(\mu^N_{t-})\}} \dif u  \dif \mathbbm{v}.
    \]
    Evaluating the integral over the auxiliary variable $u$, we invoke the identity $\int_0^\infty \mathbbm{1}_{\{u \le \Lambda(\mathbbm{z})\}} \mathrm{d}u = \Lambda(\mathbbm{z})$, yielding the continuous generator:
    \begin{equation*}
    \mathcal{L}_t f(\mathbbm{z}) = \mathcal{L}^{\text{diff}}_t f(\mathbbm{z}) + \Lambda(\mu^N_{t-}) \int_{[0,1]^N} \left( f(\Psi(\mathbbm{z}, \mathbbm{v})) - f(\mathbbm{z}) \right)  \mathrm{d}\mathbbm{v}.
    \end{equation*}
    
    \textbf{Step 2: Generator of the Discrete Process.}
    We now analyze the discrete generator 
    $$
        L^{\Delta t}_t f(\mathbbm{z}) = \Delta t^{-1} \mathbb{E}[f(\mathbbm{z}_{t+\Delta t}) - f(\mathbbm{z}) \mid \mathbbm{z}_t = \mathbbm{z}].
    $$ 
    The process evolves via diffusion with probability $1 - \Lambda(\mu^N_{t-})\Delta t$ and undergoes resampling with probability $\Lambda(\mu^N_{t-})\Delta t$.
    
    For the diffusion step, Taylor expansion implies 
    $$
        \mathbb{E}_{\text{diff}}[f(\mathbbm{z}_{t+\Delta t})] = f(\mathbbm{z}) + \mathcal{L}^{\text{diff}}_t f(\mathbbm{z})\Delta t + o(\Delta t).
    $$ 
    For the resampling step, the expectation over the multinomial selection of indices is given by:
    \[
    \mathbb{E}_{\text{res}}[f(\mathbbm{z}_{\text{new}})] = \sum_{k_1=1}^N \dots \sum_{k_N=1}^N \left( \prod_{i=1}^N \widehat{w}^{(k_i)} \right) f\left( \mathbbm{z}(k_1, \dots, k_N) \right).
    \]
    where $\mathbbm{z}(k_1, \dots, k_N)$ represents the state where the $i$-th particle is replaced by the $k_i$-th parent.
    
    We now show this discrete sum is equivalent to the integral representation used in the continuous generator. Consider the auxiliary variables $\mathbbm{v} = (v_1, \dots, v_N) \in [0,1]^N$. The mapping $\Psi(\mathbbm{z}, \mathbbm{v})$ assigns particle $i$ to parent $k$ if $v_i$ falls into the $k$-th cumulative weight interval $I_k=(\sum_{i=1}^{k-1}\widehat{w}^{(i)},\sum_{i=1}^{k}\widehat{w}^{(i)}]$, where $|I_k| = \widehat{w}^{(k)}$.
    Since $v_1, \dots, v_N$ are independent uniform random variables, the integral over the hypercube factors as:
    \begin{align*}
    \int_{[0,1]^N} f(\Psi(\mathbbm{z}, \mathbbm{v}))  \dif \mathbbm{v}
    &= \int_0^1 \dots \int_0^1 f\left( \mathbbm{z}\left(\sum_{k_1} k_1 \mathbbm{1}_{I_{k_1}}(v_1), \dots, \sum_{k_N} k_N \mathbbm{1}_{I_{k_N}}(v_N)\right) \right) dv_1 \dots dv_N \\
    &= \sum_{k_1=1}^N \dots \sum_{k_N=1}^N \left( \int_{I_{k_1}} dv_1 \dots \int_{I_{k_N}} dv_N \right) f\left( \mathbbm{z}(k_1, \dots, k_N) \right) \\
    &= \sum_{k_1=1}^N \dots \sum_{k_N=1}^N \left( \prod_{i=1}^N \widehat{w}^{(k_i)} \right) f\left( \mathbbm{z}(k_1, \dots, k_N) \right).
    \end{align*}
    This establishes the following identity:
    \[
        \mathbb{E}_{\text{res}}[f(\mathbbm{z}_{\text{new}})] = \int_{[0,1]^N} f(\Psi(\mathbbm{z}, \mathbbm{v}))  \dif \mathbbm{v}.
    \]
    
    Substituting the results back into the expression for $L^{\Delta t}_t f(\mathbbm{z})$:
    \begin{align*}
    L^{\Delta t}_t f(\mathbbm{z}) &= \frac{1}{\Delta t} \Big[ (1 - \Lambda(\mu^N_{t-})\Delta t)(f(\mathbbm{z}) + \mathcal{L}_{\text{diff}}f(\mathbbm{z})\Delta t) + \Lambda(\mu^N_{t-})\Delta t \int_{[0,1]^N} f(\Psi(\mathbbm{z}, \mathbbm{v})) \dif \mathbbm{v} - f(\mathbbm{z}) + o(\Delta t) \Big] \\
    &= \mathcal{L}^{\text{diff}}_tf(\mathbbm{z}) - \Lambda(\mu^N_{t-}) f(\mathbbm{z}) + \Lambda(\mu^N_{t-}) \int_{[0,1]^N} f(\Psi(\mathbbm{z}, \mathbbm{v})) \dif \mathbbm{v} + o(1) \\
    &= \mathcal{L}^{\text{diff}}_t f(\mathbbm{z}) + \Lambda(\mu^N_{t-}) \int_{[0,1]^N} \left( f(\Psi(\mathbbm{z}, \mathbbm{v})) - f(\mathbbm{z}) \right) \dif \mathbbm{v} + o(1).
    \end{align*}
    Taking the limit as \(\Delta t \to 0\), we obtain 
    \[
        \lim_{\Delta t \to 0} L^{\Delta t}_t f(\mathbbm{z}) = \mathcal{L}_t f(\mathbbm{z}).
    \] 
    Since the generator of the discrete approximation converges uniformly to the generator of the continuous jump-diffusion process, the discrete system converges weakly to the continuous system.
\end{proof}

\end{document}
